Para evaluar la utilidad del marco más allá de la hora siguiente, se repitió el protocolo de origen móvil con pronóstico directo a $h \in \{1, 6, 12, 24, 48\}$ horas (Tabla~\ref{tab_horizons}, Figura~\ref{fig_horizons}). Cada modelo se entrenó específicamente para cada horizonte y las líneas base utilizaron la misma información disponible en el origen. En ATE, el MAE del mejor modelo pasó de 3.12 (HGB, $h=1$) a 5.63\,$\mu$g/m$^3$ (HGB, $h=48$), un aumento de 80\%; superó al Naive en $h \in \{1, 6, 12, 24, 48\}$ y al Naive estacional de 24 h en $h \in \{1, 6, 12, 24, 48\}$. En Campo de Marte, el MAE del mejor modelo pasó de 6.10 (RF, $h=1$) a 14.28\,$\mu$g/m$^3$ (RF, $h=48$), un aumento de 134\%; superó al Naive en $h \in \{1, 6, 12\}$ y al Naive estacional de 24 h en $h \in \{1, 6, 12\}$. En San Borja, el MAE del mejor modelo pasó de 3.22 (MLP, $h=1$) a 6.86\,$\mu$g/m$^3$ (HGB, $h=48$), un aumento de 113\%; superó al Naive en $h \in \{1, 6, 12\}$ y al Naive estacional de 24 h en $h \in \{1\}$. En Santa Anita, el MAE del mejor modelo pasó de 3.36 (MLP+met, $h=1$) a 7.55\,$\mu$g/m$^3$ (SARIMA, $h=48$), un aumento de 125\%; superó al Naive en $h \in \{1, 6, 12, 48\}$ y al Naive estacional de 24 h en $h \in \{1, 48\}$. En los horizontes de 12 a 48 h, el mejor modelo superó al Naive estacional de 24 h en 5 de 12 combinaciones estación--horizonte (42\%). Este resultado delimita el alcance operativo de la formulación univariada: su ventaja se concentra en el corto plazo, mientras que a horizontes diarios el perfil horario persistente captura buena parte de la información predecible.
